\documentclass[a4paper,twoside]{article}

\usepackage{epsfig}
\usepackage{graphicx}
\usepackage{subcaption}
\usepackage{calc}
\usepackage{amssymb}
\usepackage{amstext}
\usepackage{amsmath}
\usepackage{amsthm}
\usepackage{multicol}
\usepackage{pslatex}
\usepackage{apalike}
\usepackage{algorithm2e}
\usepackage[bottom]{footmisc}
\usepackage{pifont}
\usepackage{array} 
\usepackage{bm}
\usepackage{xcolor}
\usepackage{colortbl}
\usepackage[utf8]{inputenc}
\usepackage{url}
\usepackage[hidelinks]{hyperref}
%\usepackage{remreset}
\usepackage{booktabs}
\usepackage{diagbox}
\newcolumntype{P}[1]{>{\centering\arraybackslash}m{#1}}
\usepackage{doi}
\usepackage{multirow}
\usepackage{siunitx}
%\usepackage{pdfpages}
\usepackage{threeparttable}
\usepackage{arydshln}
\usepackage{xurl}
\usepackage{SCITEPRESS}     % Please add other packages that you may need BEFORE the SCITEPRESS.sty package.

\begin{document}

\title{Adaptive RAG Configuration Selection Based on Question and Document Features Considering Cost and Accuracy}

\author{anonymous for review}
%\author{\authorname{Fumikatsu Anaguchi\sup{1,2}, Sudesna Chakraborty\sup{1}, Takeshi Morita\sup{1,2}, Aoi Ohta\sup{2}, Masaki Asada\sup{2}, Shusaku Egami\sup{2}, Takanori Ugai\sup{3,2} and Masahiro Hamasaki\sup{2}}
%\affiliation{\sup{1}Aoyama Gakuin University, Japan}
%\affiliation{\sup{2}National Institute of Advanced Industrial Science and Technology (AIST), Japan}
%\affiliation{\sup{3}Fujitsu Limited, Japan}
%\email{morita@it.aoyama.ac.jp}
%}

\keywords{Retrieval-Augmented Generation, Graph-Based RAG, Adaptive RAG, Learning to Rank, Large Language Model.}

\abstract{Large language models (LLMs) may generate misinformation or hallucinate facts that are absent from their training data. Retrieval-augmented generation (RAG) addresses this limitation by retrieving external knowledge and providing it to an LLM as context. However, different queries may require different RAG configurations: complex queries can benefit from high-cost strategies (e.g., GraphRAG), whereas simple queries may be answered using fast, low-cost configurations (e.g., NaiveRAG). Adaptive RAG systems therefore aim to balance response quality and computational cost. Existing approaches often rely on LLM-based query routing, which introduces runtime overhead, or online reinforcement learning, which incurs ongoing operational costs. Moreover, few unified frameworks optimize multiple objectives, such as accuracy and latency, across heterogeneous configurations that vary in chunk size, embedding model, and generator model. This study proposes an offline learning-to-rank approach for adaptive RAG configuration selection. We train a ranking model on query and document features to rank 96 candidate RAG configurations. At inference time, the model ranks candidates without additional LLM calls and selects a configuration according to user-defined priority weights for precision, recall, $F_1$ score, and execution time. Evaluation results show that our method outperforms standard baselines and achieves Pareto improvements. Specifically, the balanced configuration improves BERTScore $F_1$ by 14.15 percentage points and reduces execution time by 86.2\% relative to the Adaptive-RAG baseline. The reported Wilcoxon signed-rank tests show statistically significant improvements over Adaptive-RAG.}

\onecolumn \maketitle \normalsize \setcounter{footnote}{0} \vfill

\input{ICAART2027/sec1}
\input{ICAART2027/sec2}
\input{ICAART2027/sec3}
\input{ICAART2027/sec4}
\input{ICAART2027/sec5}
\input{ICAART2027/sec6}

%\section*{\uppercase{Acknowledgements}}
%This work was supported in part by New Energy and Industrial Technology Development Organization (NEDO) under Grant Numbers JPNP20006 and JPNP25006, and by JSPS KAKENHI Grant Numbers 23K11221 and 25K03232.

\bibliographystyle{apalike}
{\small
\bibliography{references}}

\end{document}
